\chapter{Experimental Results and Analysis}
\label{chap:results}

This chapter reports the empirical results obtained with the task formulations
defined in \cref{chap:task-formulations} and the experimental protocol defined
in \cref{chap:experimental-setup}. The analysis follows the same order as the
modeling pipeline: first the prepared data are summarized, then each task is
evaluated in its native predictive space, and finally all model outputs are
converted into switch decisions and compared against the common Perfect Switch
reference.

The test set is always the external file \texttt{fc-uplink-fade.csv}. Native
metrics are useful for diagnosing whether a model learned its assigned target,
but the main operational interpretation is based on the switch metrics, because
the practical objective is the decision of whether the backup link should be
activated.

Throughout this chapter, arrows in performance-table headers indicate the
preferred direction of each metric: \(\uparrow\) means that higher values are
better, whereas \(\downarrow\) means that lower values are better. Descriptive
quantities without a universally preferred direction are shown without an
arrow.

\section{Preliminary Data Analysis}
\label{sec:results-preliminary-data-analysis}

The raw collection contains eight signal files acquired between July 2020 and
July 2021. After parsing timestamps, selecting the signal column, removing
invalid rows, and enforcing the signal-only schema, the usable collection
contains 2,835,814 valid observations out of 3,042,803 raw rows. Four files
contain additional columns: the weather-beacon file contains meteorological
variables, the two \texttt{fc-*} files contain one extra column, and
\texttt{uplink\_beacon\_data\_fucino\_GX4.csv} contains three additional
columns. These fields are ignored in all experiments, consistently with the
signal-only policy.

The sampling analysis confirms that the data are close to a regular
\SI{30}{s} grid. Six files have median inter-sample spacing of \SI{30}{s}, while
the GX4 uplink beacon file has median spacing of \SI{31}{s}. The complete
collection contains 462 continuous acquisition
segments when gaps larger than \SI{90}{s} are used as segment boundaries. These
boundaries are important because input windows and duration targets must not
assume continuity across long gaps. Short gaps are handled only by the
conservative imputation policy described in
\cref{sec:experimental-dataset-configuration}.

Most samples are close to the non-fade operating region, while threshold
exceedances are rare and concentrated in events. For example, the median signal
is between approximately \SI{0.3}{dB} and \SI{1.5}{dB} across the source files,
whereas the final operational threshold is \(\theta=\SI{10.0}{dB}\). The
largest observed values are much higher, reaching \SI{69.0}{dB} in the
Fucino uplink sources. This imbalance is the reason why all tasks are built
around candidate fades, grouped fade events, or residual durations rather than
around uniformly sampled background windows.

\begin{table}[htbp]
  \centering
  \caption{Final supervised datasets used in the experimental phase.}
  \label{tab:final-supervised-datasets}
  \small
  \setlength{\tabcolsep}{4pt}
  \renewcommand{\arraystretch}{1.2}
  \begin{tabular}{
    @{}
    >{\raggedright\arraybackslash}p{0.25\textwidth}
    >{\raggedright\arraybackslash}p{0.18\textwidth}
    rrr
    >{\raggedright\arraybackslash}p{0.21\textwidth}
    @{}
  }
    \toprule
    Task & Main target & Train & Validation & Test & Event coverage \\
    \midrule
    Autoregressive forecasting
      & future signal trajectory
      & 19,412 & 4,783 & 7,708
      & 61/15/24 events \\
    Current-level persistence
      & residual persistence above \(S_i-\delta\)
      & 19,598 & 4,798 & 7,692
      & 61/15/24 events \\
    Long-fade detection
      & long-event class label
      & 4,589 & 790 & 2,569
      & 62/15/25 grouped events \\
    Survival persistence
      & residual fade-event duration
      & 1,729 & 729 & 921
      & 117/23/55 survival events \\
    \bottomrule
  \end{tabular}
\end{table}

\Cref{tab:final-supervised-datasets} shows that the four task datasets do not
have the same number of samples, even though they originate from the same raw
signals and use the same external holdout source. This is expected. The
autoregressive and current-level datasets are built from candidate event
windows; long-fade detection is restricted to timestamps inside grouped fade
events; survival persistence is restricted to timestamps inside state-machine
fade events and keeps censored observations. Consequently, cross-task
comparison cannot rely only on native losses. The common switch-conversion and
reference-grid analyses later in this chapter are used to compare the
operational behavior of models trained under these different sample
definitions.

\section{Autoregressive Forecasting}
\label{sec:results-autoregressive}

The autoregressive experiments predict the next \(h=10\) signal samples from a
past context of \(L=30\) samples. Validation RMSE is reported for model
selection; final forecast and switch metrics use the external test file.
Forecast errors on both splits are reported in raw signal scale, while switch
metrics are test-only.
The switch decision is obtained from the predicted trajectory: a switch is
requested at time \(i\) only when all 10 predicted future samples are at least
\(\SI{10.0}{dB}\), followed by the shared switch post-processing defined in
\cref{sec:shared-switch-post-processing}.

\subsection{Forecast Accuracy}
\label{sec:results-autoregressive-forecast-accuracy}

\Cref{tab:autoregressive-forecast-results} summarizes validation and test
forecast errors. The validation column is absent for Chronos because it is used
only in zero-shot mode and does not load train or validation data.

\begin{table}[htbp]
  \centering
  \caption{Autoregressive validation RMSE used for model selection and forecast
  accuracy on the external holdout split.}
  \label{tab:autoregressive-forecast-results}
  \footnotesize
  \setlength{\tabcolsep}{3pt}
  \renewcommand{\arraystretch}{1.15}
  \begin{tabular}{
    @{}
    >{\raggedright\arraybackslash}p{0.35\textwidth}
    rrr
    >{\raggedright\arraybackslash}p{0.16\textwidth}
    @{}
  }
    \toprule
    Method & Val. RMSE \(\downarrow\) & Test MAE \(\downarrow\) & Test RMSE \(\downarrow\) & Variant \\
    \midrule
    GRU S2V & 3.124 & 1.343 & 2.852 & raw \\
    GRU Seq2Seq & 3.560 & 1.326 & 3.040 & raw \\
    XGBoost & 3.863 & 1.376 & 3.088 & raw \\
    PatchTST & 4.090 & 1.515 & 3.088 & raw \\
    XGBoost & 3.221 & 1.465 & 3.363 & context std. \\
    PatchTST & 3.811 & 1.475 & 3.512 & context std. \\
    GRU S2V & 3.686 & 1.443 & 3.558 & context std. \\
    GRU Seq2Seq & 3.702 & 1.460 & 3.671 & context std. \\
    Chronos zero-shot & -- & 1.540 & 3.899 & raw \\
    \bottomrule
  \end{tabular}
\end{table}

The best test RMSE is obtained by the raw GRU sequence-to-vector model
(\(2.852\)), while the lowest test MAE is obtained by the raw GRU
encoder-decoder (\(1.326\)). This difference is expected because RMSE penalizes
large peak errors more strongly than MAE. The raw XGBoost and raw PatchTST
models have almost identical test RMSE, despite using very different model
classes. Chronos performs worse than the trained models on this external
holdout, which is consistent with its zero-shot use: it is not adapted to the
particular scale, noise, and event morphology of the Fucino data.

The context-standardized variants do not improve external test performance.
Although context standardization can simplify optimization, it removes part of
the absolute signal-level information that is operationally important near the
\SI{10.0}{dB} threshold. For this task, the raw representation is therefore the
more reliable choice.

\Cref{fig:autoregressive-patchtst-horizon-error} shows the horizon-wise error
profile for PatchTST raw. Both MAE and RMSE generally increase with the forecast step,
confirming that the last points of the \SI{300}{s} horizon are more uncertain
than the immediate future. This behavior matters for switch conversion because
the operational rule requires all 10 predicted samples to remain above the
threshold.

\begin{figure}[H]
  \centering
  \includegraphics[width=0.82\textwidth]{results/autoregressive_patchtst_raw_horizon_error.png}
  \caption{Horizon-wise forecast error for the PatchTST raw autoregressive model on the external holdout split.}
  \label{fig:autoregressive-patchtst-horizon-error}
\end{figure}

\subsection{Switch Performance}
\label{sec:results-autoregressive-switch-performance}

\Cref{tab:autoregressive-switch-results} reports the switch metrics obtained
after applying the trajectory-to-switch rule and the shared post-processing.
The ranking differs from the pure forecast-error ranking. PatchTST raw has the
best F1 and IoU, even though it is not the best model in RMSE.

\begin{table}[htbp]
  \centering
  \caption{Autoregressive switch performance against Perfect Switch on the external holdout split.}
  \label{tab:autoregressive-switch-results}
  \small
  \setlength{\tabcolsep}{5pt}
  \renewcommand{\arraystretch}{1.15}
  \begin{tabular}{
    @{}
    >{\raggedright\arraybackslash}p{0.38\textwidth}
    rrrr
    @{}
  }
    \toprule
    Method & Precision \(\uparrow\) & Recall \(\uparrow\) & F1 \(\uparrow\) & IoU \(\uparrow\) \\
    \midrule
    PatchTST raw & 0.779 & 0.839 & 0.808 & 0.678 \\
    XGBoost raw & 0.692 & 0.943 & 0.798 & 0.664 \\
    GRU Seq2Seq raw & 0.619 & 0.980 & 0.759 & 0.611 \\
    Chronos zero-shot raw & 0.567 & 0.961 & 0.713 & 0.554 \\
    XGBoost context-standard & 0.512 & 0.991 & 0.675 & 0.510 \\
    GRU S2V raw & 0.509 & 0.993 & 0.674 & 0.508 \\
    PatchTST context-standard & 0.496 & 0.991 & 0.661 & 0.493 \\
    GRU S2V context-standard & 0.496 & 0.987 & 0.660 & 0.493 \\
    GRU Seq2Seq context-standard & 0.497 & 0.985 & 0.660 & 0.493 \\
    \bottomrule
  \end{tabular}
\end{table}

Point-wise precision, recall, F1, and IoU do not fully describe the operational
behavior of a switch policy. Two models can have similar F1 while producing
different active durations or a different number of switch activations. For this
reason, \cref{fig:autoregressive-switch-duration-events} reports the total
post-processed switch duration and the number of contiguous switch islands on
the external holdout split. The Perfect Switch contains 459 active samples,
corresponding to \SI{229.5}{min}, grouped into 18 switch islands. PatchTST raw
is the closest autoregressive model in this diagnostic view, with 494 active
samples and 20 islands. Several other models reach high recall by keeping the
switch active for much longer or by producing many more islands, which explains
their lower precision.

\begin{figure}[H]
  \centering
  \includegraphics[width=0.8\textwidth]{results/autoregressive_switch_duration_events.png}
  \caption{Autoregressive switch duration and event-count behavior on the external holdout split. Duration counts post-processed positive switch samples, converted to minutes using the \SI{30}{s} sampling time. The dashed red line marks the Perfect Switch value.}
  \label{fig:autoregressive-switch-duration-events}
\end{figure}

The common pattern is high recall and lower precision. In operational terms,
most Perfect Switch points are detected, but several models keep the switch
active for longer than the reference. The raw PatchTST model is the most
balanced among the autoregressive models: it sacrifices some recall relative to
the GRU models, but reduces false positives enough to obtain the highest F1 and
IoU. XGBoost raw is close behind and is notable because it reaches competitive
switch performance without an explicit neural sequence architecture.

The confusion counts clarify this trade-off. PatchTST raw detects 385 of 459
Perfect Switch samples and misses 74, while adding 109 false-positive switch
samples over 7708 evaluated timestamps. This corresponds to a balanced accuracy
of 0.912 and a false-positive rate of 0.015. Its advantage is therefore not
that it detects every switch sample, but that it preserves most operational
intervals while keeping unnecessary switch time limited. The
context-standardized variants have the opposite tendency: they often increase
recall, but they lose precision because local standardization removes part of
the absolute-level information needed near the \SI{10.0}{dB} threshold.

\subsection{Interpretation and Event-Level Examples}
\label{sec:results-autoregressive-examples}


\begin{figure}[H]
  \centering
  \includegraphics[width=\textwidth]{results/autoregressive_patchtst_raw_event_00007_switch_comparison.png}
  \caption{Example autoregressive switch comparison for PatchTST raw on a multi-interval external-holdout event. The model-derived switch is compared with the Perfect Switch reference after applying the common post-processing rule.}
  \label{fig:autoregressive-patchtst-switch-example}
\end{figure}

\Cref{fig:autoregressive-patchtst-switch-example} shows event
\texttt{dataset\_007\_event\_00007}, which is more informative than a single
isolated fade because the Perfect Switch turns on and off several times. The
upper panel contains the Perfect Switch reference, while the lower panel
contains the PatchTST raw decision. The model captures the long central switch
interval and the late reactivation, but misses the first short
activation. This illustrates why switch metrics are not determined only by
point forecast error: a model can tolerate moderate forecast deviations if it
still preserves the operational threshold crossing structure.


\Cref{fig:autoregressive-switch-mosaic} compares all autoregressive switch
decisions on event 00003 of dataset 007, a different, single-interval example
from \cref{fig:autoregressive-patchtst-switch-example}.
The upper panel shows the raw signal and the
\SI{10.0}{dB} threshold, while each lower row shows the final post-processed
switch sequence for one method. This compact view makes clear that most
autoregressive models identify the main fade, but they slightly differ in onset timing,
release timing, and the amount of extra switch activity around the event.

\begin{figure}[H]
  \centering
  \includegraphics[width=0.95\textwidth]{results/autoregressive_event00003_switch_mosaic.png}
  \caption{Mosaic of post-processed switch decisions for all autoregressive methods on one external-holdout event. Each row reports the final switch used for evaluation, not the raw thresholded model output.}
  \label{fig:autoregressive-switch-mosaic}
\end{figure}

\section{Current-Level Persistence}
\label{sec:results-current-level-persistence}

The current-level persistence task predicts how long the signal will remain
above \(S_i-\delta\), with \(\delta=\SI{0.5}{dB}\). This is a scalar duration
regression problem, not a future-trajectory forecast. A switch is requested only
when the current signal is at least \(\SI{10.0}{dB}\) and the predicted
remaining persistence is at least \SI{300}{s}; the same switch post-processing
used for the other tasks is then applied.

\subsection{Duration Prediction Accuracy}
\label{sec:results-current-level-duration}

\Cref{tab:current-level-duration-results} reports the validation and test
duration-regression metrics. Errors are expressed in seconds because the
operational threshold is \SI{300}{s}. The table also reports the median
absolute error, which is important because the target distribution is highly
right-skewed: a small number of very long persistence intervals can dominate
MAE and especially RMSE.

\begin{table}[htbp]
  \centering
  \caption{Current-level persistence duration-regression results on the external holdout split. All error metrics are expressed in seconds.}
  \label{tab:current-level-duration-results}
  \small
  \setlength{\tabcolsep}{4pt}
  \renewcommand{\arraystretch}{1.15}
  \begin{tabular}{
    @{}
    >{\raggedright\arraybackslash}p{0.34\textwidth}
    rrrr
    @{}
  }
    \toprule
    Method
      & \shortstack{Val. MAE\\\((\mathrm{s})\,\downarrow\)}
      & \shortstack{Test MAE\\\((\mathrm{s})\,\downarrow\)}
      & \shortstack{Test RMSE\\\((\mathrm{s})\,\downarrow\)}
      & \shortstack{Test median AE\\\((\mathrm{s})\,\downarrow\)} \\
    \midrule
    XGBoost scalar context & 1220.9 & 586.9 & 2584.6 & 27.2 \\
    Shapelet convolution scalar context & 4613.2 & 1835.6 & 6491.0 & 141.5 \\
    Shapelet MLP scalar context & 5116.5 & 2063.8 & 7629.7 & 160.6 \\
    Shapelet Transformer scalar context & 5278.9 & 2056.6 & 7685.3 & 169.4 \\
    \bottomrule
  \end{tabular}
\end{table}

XGBoost is clearly the strongest model for this formulation. Its test median
absolute error is only \SI{27.2}{s}, which means that typical short-horizon
decisions are often accurate. However, its test MAE is \SI{586.9}{s} and its
RMSE is \SI{2584.6}{s}, showing that large errors still occur on long
persistence intervals. The learnable-shapelet models are substantially worse:
their test MAE is around \SI{30}{min} to \SI{34}{min}, and their RMSE exceeds
\SI{108}{min}. These errors are too large for the scalar duration prediction to
be considered reliable by itself.

\Cref{fig:current-level-duration-errors} summarizes the same behavior in
minutes. The median errors are much smaller than the mean errors, confirming
that the hard cases are concentrated in the long tail of the duration
distribution rather than uniformly distributed over all windows.

\begin{figure}[H]
  \centering
  \includegraphics[width=\textwidth]{results/current_level_duration_errors.png}
  \caption{Current-level persistence duration-regression errors on the external holdout split. MAE and median absolute error describe typical behavior, while RMSE highlights sensitivity to long-duration mistakes.}
  \label{fig:current-level-duration-errors}
\end{figure}

\Cref{fig:current-level-xgboost-duration-scatter} shows the predicted and true
durations for the best current-level model. The dense region near the origin is
well aligned, but the dispersion increases for longer intervals. This explains
why the median absolute error is low while MAE and RMSE remain large.

\begin{figure}[H]
  \centering
  \includegraphics[width=0.60\textwidth]{results/current_level_xgboost_duration_predicted_vs_true.png}
  \caption{Predicted versus true current-level persistence duration for the XGBoost scalar-context model on the external holdout split.}
  \label{fig:current-level-xgboost-duration-scatter}
\end{figure}

\subsection{Switch Performance}
\label{sec:results-current-level-switch}

\Cref{tab:current-level-switch-results} reports the resulting switch metrics.
Despite the imperfect duration regression, XGBoost gives a useful switch
policy: it reaches an F1 score of \(0.902\) and produces no false-positive
switch samples on the external holdout split. The three learnable-shapelet
models instead produce an all-zero final switch sequence, so they obtain zero
recall and zero F1. This is not caused by missing predictions. The signal gate
is active at 672 test timestamps, but the maximum persistence predicted at
those timestamps is only \SI{148.8}{s} for the shapelet MLP, \SI{204.6}{s} for
the shapelet Transformer, and \SI{167.3}{s} for the shapelet convolution model.
All remain below the \SI{300}{s} duration threshold. The same models do predict
durations above \SI{300}{s} at other timestamps, but only when the current
signal is below \SI{10.0}{dB}. Consequently, the two conditions required by the
decision rule are never true simultaneously. This is an operationally
important negative result: the models fit the regression target to some extent,
but they are not usable switch controllers under the configured rule.

\begin{table}[htbp]
  \centering
  \caption{Current-level persistence switch performance against Perfect Switch on the external holdout split.}
  \label{tab:current-level-switch-results}
  \small
  \setlength{\tabcolsep}{5pt}
  \renewcommand{\arraystretch}{1.15}
  \begin{tabular}{
    @{}
    >{\raggedright\arraybackslash}p{0.38\textwidth}
    rrrrr
    @{}
  }
    \toprule
    Method & Precision \(\uparrow\) & Recall \(\uparrow\) & F1 \(\uparrow\) & IoU \(\uparrow\) & Active min. \\
    \midrule
    XGBoost scalar context & 1.000 & 0.822 & 0.902 & 0.822 & 189.5 \\
    Shapelet MLP scalar context & 0.000 & 0.000 & 0.000 & 0.000 & 0.0 \\
    Shapelet Transformer scalar context & 0.000 & 0.000 & 0.000 & 0.000 & 0.0 \\
    Shapelet convolution scalar context & 0.000 & 0.000 & 0.000 & 0.000 & 0.0 \\
    \bottomrule
  \end{tabular}
\end{table}

\Cref{fig:current-level-switch-duration-events} makes this contrast explicit.
The Perfect Switch has 461 active samples, corresponding to \SI{230.5}{min},
and 17 switch islands. XGBoost produces 379 active samples and 11 islands. It
therefore under-covers part of the Perfect Switch, but its conservative
behavior avoids the excessive switch activity observed in several
autoregressive variants. The shapelet models produce no active switch interval
under the configured decision rule.

This result should be read together with the decision rule. A current-level
model is allowed to request a switch only when the current signal is already at
or above the \SI{10.0}{dB} threshold and the predicted persistence exceeds
\SI{300}{s}. The XGBoost model therefore behaves as a conservative switch
policy: it produces 379 true-positive switch samples, no false-positive samples,
and 82 false negatives. The resulting precision of 1.000 is operationally
attractive because every requested switch is aligned with the reference, but
the recall loss indicates that some useful backup-link intervals would be
missed. This is a different failure mode from autoregressive forecasting, where
the main issue is excess active duration rather than missed coverage.

\begin{figure}[H]
  \centering
  \includegraphics[width=\textwidth]{results/current_level_switch_duration_events.png}
  \caption{Current-level persistence switch duration and event-count behavior on the external holdout split. Duration counts post-processed positive switch samples, converted to minutes using the \SI{30}{s} sampling time.}
  \label{fig:current-level-switch-duration-events}
\end{figure}

\subsection{Interpretation and Event-Level Examples}
\label{sec:results-current-level-examples}


\begin{figure}[H]
  \centering
  \includegraphics[width=0.9\textwidth]{results/current_level_xgboost_event_00007_switch_comparison.png}
  \caption{Example current-level persistence switch comparison for XGBoost scalar context on the selected multi-interval external-holdout event. The model-derived switch is compared with the Perfect Switch reference after applying the common post-processing rule.}
  \label{fig:current-level-xgboost-switch-example}
\end{figure}



\Cref{fig:current-level-xgboost-switch-example} shows the current-level
XGBoost model on the task-native decision grid for the same source event
identifier used in \cref{fig:autoregressive-patchtst-switch-example}. Because each task has its own
valid decision timestamps, the plotted Perfect Switch can differ slightly
between task-native figures. The model decision follows the long central
interval well and avoids the short early activation. This is consistent with
the global behavior of the task: the model is conservative and does not create
false-positive switch samples, but it can ignore short or marginal intervals
that the Perfect Switch keeps. The example also clarifies the role of the
formulation: the model does not forecast the whole future trajectory; it only
estimates whether the current level is likely to persist long enough to justify
the switch.


\section{Long-Fade Detection}
\label{sec:results-long-fade-detection}

The long-fade detection task is a binary classification problem defined inside
grouped fade events. At each eligible timestamp, the model estimates whether
the grouped event will last at least \SI{300}{s}. The native output is therefore
a class probability, not a duration or a future signal trajectory. For switch
evaluation, a switch is requested when the predicted long-fade probability is at
least \(0.5\), followed by the same shared switch post-processing used in the
other tasks.

\subsection{Classification Metrics}
\label{sec:results-long-fade-classification}

\Cref{tab:long-fade-classification-results} reports the native classification
metrics. All three models reach very high validation and test AUPRC. This
indicates that the long-fade label is easy to separate once the sample is known
to belong to a grouped fade event. The result must therefore be interpreted with
care: the task is useful as an event-level persistence classifier, but its
native metrics are less demanding than the final switch decision because the
dataset excludes most non-event background timestamps.

\begin{table}[htbp]
  \centering
  \caption{Long-fade detection classification results on the external holdout split.}
  \label{tab:long-fade-classification-results}
  \small
  \setlength{\tabcolsep}{4pt}
  \renewcommand{\arraystretch}{1.15}
  \begin{tabular}{
    @{}
    >{\raggedright\arraybackslash}p{0.36\textwidth}
    rrrr
    @{}
  }
    \toprule
    Method
      & \shortstack{Val. AUPRC\\\(\uparrow\)}
      & \shortstack{Test AUPRC\\\(\uparrow\)}
      & \shortstack{Test AUROC\\\(\uparrow\)}
      & \shortstack{Test F1@0.5\\\(\uparrow\)} \\
    \midrule
    TCN classifier & 1.000 & 1.000 & 1.000 & 0.999 \\
    XGBoost lag + scalar & 1.000 & 1.000 & 1.000 & 1.000 \\
    Shapelet convolution classifier & 1.000 & 1.000 & 0.999 & 0.997 \\
    \bottomrule
  \end{tabular}
\end{table}

The TCN and XGBoost classifiers are almost indistinguishable in native
classification space. The shapelet convolution model is slightly weaker on the
test split, but it remains close to the other two models. The stronger
differences appear only after the probabilities are converted into switch
sequences and compared with the Perfect Switch reference.

\subsection{Switch Performance}
\label{sec:results-long-fade-switch}

\Cref{tab:long-fade-switch-results} reports the switch metrics obtained after
thresholding the predicted probabilities and applying the shared
post-processing. All models recover nearly all Perfect Switch samples, but this
comes at the cost of many false-positive switch samples. The TCN has the best
F1 and IoU, closely followed by XGBoost.

\begin{table}[htbp]
  \centering
  \caption{Long-fade detection switch performance against Perfect Switch on the external holdout split.}
  \label{tab:long-fade-switch-results}
  \small
  \setlength{\tabcolsep}{5pt}
  \renewcommand{\arraystretch}{1.15}
  \begin{tabular}{
    @{}
    >{\raggedright\arraybackslash}p{0.38\textwidth}
    rrrrr
    @{}
  }
    \toprule
    Method & Precision \(\uparrow\) & Recall \(\uparrow\) & F1 \(\uparrow\) & IoU \(\uparrow\) & Active min. \\
    \midrule
    TCN classifier & 0.540 & 1.000 & 0.702 & 0.540 & 445.0 \\
    XGBoost lag + scalar & 0.540 & 1.000 & 0.701 & 0.540 & 445.5 \\
    Shapelet convolution classifier & 0.521 & 0.998 & 0.685 & 0.521 & 460.5 \\
    \bottomrule
  \end{tabular}
\end{table}

\Cref{fig:long-fade-switch-duration-events} shows why the switch scores are
lower than the native classification scores. The Perfect Switch contains 481
active samples, corresponding to \SI{240.5}{min}, and 19 switch islands on the
long-fade evaluation grid. The model-derived switches contain almost twice as
many active samples and more than twice as many islands. This means that the
models correctly identify long fades, but they keep the switch active over a
broader portion of the grouped-event windows than the Perfect Switch requires.

For the TCN classifier, the confusion matrix contains no false negatives: all
481 Perfect Switch samples are detected. The cost is 409 false-positive
samples, a false-positive rate of 0.196, and 890 active switch samples in
total. This explains why the recall is 1.000 but the precision is only 0.540.
The result is not a modeling failure in the native classification task; it is
a consequence of converting a grouped-event label into an instantaneous switch
policy. Once a grouped fade is classified as long, the model tends to remain
positive over a wide portion of that grouped event, including regions where the
Perfect Switch has already turned off.

\begin{figure}[H]
  \centering
  \includegraphics[width=\textwidth]{results/long_fade_switch_duration_events.png}
  \caption{Long-fade detection switch duration and event-count behavior on the external holdout split. Duration counts post-processed positive switch samples, converted to minutes using the \SI{30}{s} sampling time. The dashed red line marks the Perfect Switch value.}
  \label{fig:long-fade-switch-duration-events}
\end{figure}

\subsection{Interpretation and Event-Level Examples}
\label{sec:results-long-fade-examples}

\begin{figure}[H]
  \centering
  \includegraphics[width=\textwidth]{results/long_fade_tcn_event_00007_switch_comparison.png}
  \caption{Example long-fade detection switch comparison for the TCN classifier on the selected multi-interval external-holdout event. The model-derived switch is compared with the Perfect Switch reference after applying the common post-processing rule.}
  \label{fig:long-fade-tcn-switch-example}
\end{figure}

\Cref{fig:long-fade-tcn-switch-example} shows a task-native long-fade event for
the TCN classifier. The event number is local to the long-fade dataset and
should not be interpreted as the same global event number used in the previous
task-native figures. The model recognizes the grouped event as a long fade and
keeps the switch active over most of the grouped region. This makes the
formulation robust to temporary oscillations around the threshold, but it is
less precise than the Perfect Switch when the operational reference contains
several separate on/off intervals. The example therefore illustrates the main
trade-off of long-fade detection: it is a reliable high-recall detector of
problematic events, but it is too coarse to be the most selective switch
controller. The wider signal window provides context; switch values outside
the native event segment are displayed as zero and are not evaluated.



Overall, long-fade detection is highly effective at recognizing whether a
grouped fade is long, but it is less precise as a direct switch policy. It
therefore behaves as a high-recall formulation: useful when missed long fades
are especially costly, but less attractive when unnecessary switch duration is
penalized heavily.

\section{Survival Persistence}
\label{sec:results-survival-persistence}

The survival persistence task estimates the distribution of the remaining fade
duration rather than a single deterministic target. The evaluated models output
survival probabilities \(S(u \mid X_i)\), where \(u\) is a future duration
horizon. The operational switch rule uses the \SI{300}{s} horizon: a switch is
requested when \(S(\SI{300}{s}\mid X_i)\geq 0.65\) and the current signal is at
least \(\SI{10.0}{dB}\). The probability threshold \(0.65\) is retained as an
exploratory operating point from the switch-analysis stage; it should not be
interpreted as a fully validation-calibrated threshold.

\subsection{Survival Metrics}
\label{sec:results-survival-metrics}

\Cref{tab:survival-results} reports the native survival metrics on the external
holdout split. The two models are evaluated on the same censored survival
dataset. The XGBoost-AFT model uses continuous lower and upper survival bounds,
whereas the discrete-time TCN predicts hazards on a fixed duration grid. For
this reason, likelihood values are model-specific diagnostics; the comparison
below focuses on Harrell's C-index, integrated Brier score, and horizon-specific
Brier and calibration errors.

\begin{table}[htbp]
  \centering
  \caption{Survival persistence metrics on the external holdout split.}
  \label{tab:survival-results}
  \small
  \setlength{\tabcolsep}{4pt}
  \renewcommand{\arraystretch}{1.15}
  \begin{tabular}{
    @{}
    >{\raggedright\arraybackslash}p{0.34\textwidth}
    rrrr
    @{}
  }
    \toprule
    Method & C-index \(\uparrow\) & IBS \(\downarrow\) & Brier@300s \(\downarrow\) & Calib. error@300s \(\downarrow\) \\
    \midrule
    XGBoost-AFT & 0.672 & 0.117 & 0.218 & 0.123 \\
    Discrete-time TCN & 0.658 & 0.103 & 0.214 & 0.105 \\
    \bottomrule
  \end{tabular}
\end{table}

The two models show different strengths. The C-index point estimate of
XGBoost-AFT exceeds that of the discrete-time TCN by (0.014). To assess the
robustness of this difference while preserving dependence among timestamps of
the same event, a paired event-level bootstrap resampled the 55 external-holdout
events 2,000 times. The resulting 95\% percentile interval for the C-index
difference, XGBoost-AFT minus TCN, was ([-0.021, 0.049]). Because this interval
includes zero, the observed gap does not provide evidence of a statistically
significant ranking advantage and should be interpreted only as a small
numerical difference on this test set. Conversely, the discrete-time TCN has
the lower integrated Brier score and lower \SI{300}{s} calibration error,
indicating better probabilistic accuracy over the evaluated time grid. Neither
result alone determines switch quality, because the operational decision uses
only one point on the survival curve and then applies binary post-processing.

\Cref{fig:survival-tcn-curves} illustrates the TCN output using six test
windows selected near the 5th, 23rd, 41st, 59th, 77th, and 95th percentiles of
the predicted \(S(\SI{300}{s}\mid X_i)\) distribution. This replaces a purely
chronological selection, which produced several almost indistinguishable curves
from adjacent timestamps. The legend reports, for each window, its selection
quantile, predicted survival probability at the operational horizon, predicted
median residual duration, and observed residual duration. The selected
\(S(\SI{300}{s}\mid X_i)\) values range from (0.39) to (0.81), and therefore
show materially different operational predictions around the (0.65) switch
threshold. These curves are deterministic predictions for different windows,
not repeated estimates of a common population curve; consequently, a pairwise
statistical-significance test between them is not defined. Residual overlap
instead indicates that the TCN assigns similar hazard profiles to some inputs,
especially at long horizons where all predicted survival probabilities approach
zero.

\begin{figure}[H]
  \centering
  \includegraphics[width=0.98\textwidth]{results/survival_tcn_selected_survival_curves.png}
  \caption{Representative survival curves produced by the discrete-time TCN on the external holdout split. Curves are selected by quantiles of \(S(\SI{300}{s}\mid X_i)\); colors identify the windows described in the legend, and the dotted line marks the \SI{300}{s} operational horizon.}
  \label{fig:survival-tcn-curves}
\end{figure}

\subsection{Switch Performance}
\label{sec:results-survival-switch}

\Cref{tab:survival-switch-results} reports the switch metrics after converting
survival probabilities at \SI{300}{s} into binary decisions. XGBoost-AFT is the
stronger switch model: it reaches higher precision, F1, and IoU than the
discrete-time TCN. The TCN has higher recall, but it does so by keeping the
switch active much more often.

\begin{table}[htbp]
  \centering
  \caption{Survival persistence switch performance against Perfect Switch on the external holdout split.}
  \label{tab:survival-switch-results}
  \small
  \setlength{\tabcolsep}{5pt}
  \renewcommand{\arraystretch}{1.15}
  \begin{tabular}{
    @{}
    >{\raggedright\arraybackslash}p{0.38\textwidth}
    rrrrr
    @{}
  }
    \toprule
    Method & Precision \(\uparrow\) & Recall \(\uparrow\) & F1 \(\uparrow\) & IoU \(\uparrow\) & Active min. \\
    \midrule
    XGBoost-AFT & 0.828 & 0.902 & 0.863 & 0.759 & 299.0 \\
    Discrete-time TCN & 0.657 & 0.971 & 0.784 & 0.644 & 405.5 \\
    \bottomrule
  \end{tabular}
\end{table}

\Cref{fig:survival-switch-duration-events} shows the same difference from an
operational perspective. The Perfect Switch contains 549 active samples,
corresponding to \SI{274.5}{min}, and 22 switch islands on the survival
evaluation grid. XGBoost-AFT is close in duration but produces more switch
islands, while the TCN overextends the active duration and still produces more
islands than the reference.

The XGBoost-AFT switch has 495 true positives, 103 false positives, and 54
false negatives. Compared with current-level persistence, it is less
conservative: it accepts some unnecessary switch time in order to recover more
of the Perfect Switch. Compared with long-fade detection, it is much more
selective: the active duration is close to the reference, and the F1 score rises
to 0.863. This supports the methodological motivation for survival
persistence. Estimating \(P(R_i>u\mid X_i)\) provides a quantity that is already
aligned with the operational question, rather than requiring a trajectory or
event-classification output to be indirectly converted into a switch.

\begin{figure}[H]
  \centering
  \includegraphics[width=0.82\textwidth]{results/survival_switch_duration_events.png}
  \caption{Survival persistence switch duration and event-count behavior on the external holdout split. Duration counts post-processed positive switch samples, converted to minutes using the \SI{30}{s} sampling time. The dashed red line marks the Perfect Switch value.}
  \label{fig:survival-switch-duration-events}
\end{figure}

\subsection{Interpretation and Event-Level Examples}
\label{sec:results-survival-examples}

\Cref{fig:survival-xgboost-switch-example} shows a false-positive example for
XGBoost-AFT on the task-native event \texttt{sp\_event\_00007}. This identifier
is local to the survival dataset and does not denote the event used in
\cref{fig:autoregressive-patchtst-switch-example,fig:current-level-xgboost-switch-example}.
On the 18 native decision timestamps, Perfect Switch is always zero, whereas
the final model switch is active for 12 samples. The example therefore
illustrates an unnecessary activation, not a successfully recovered reference
interval. The wider signal window also contains a later peak outside this
native event segment; the displayed switches are padded with zeros there and
must not be interpreted as predictions for that peak. This distinction matters
when interpreting task-native plots: only the native timestamps contribute to
the event metrics, while \cref{fig:cross-task-event-timeline} provides the
aligned cross-task view.

\begin{figure}[H]
  \centering
  \includegraphics[width=0.92\textwidth]{results/survival_xgboost_aft_event_00007_switch_comparison.png}
  \caption[False-positive survival switch example for XGBoost-AFT.]{False-positive survival switch example for XGBoost-AFT on \texttt{sp\_event\_00007}. The model uses \(S(\SI{300}{s}\mid X_i)\geq 0.65\), the current-signal gate, and common post-processing. Perfect Switch remains zero on the native event; switches outside that segment are displayed as zero for context only.}
  \label{fig:survival-xgboost-switch-example}
\end{figure}

Overall, survival persistence is one of the most operationally meaningful
formulations because its output is a residual-duration distribution. In the
current experiments, XGBoost-AFT is the best survival switch model, while the
TCN provides better integrated probabilistic Brier performance but a less
precise binary switch sequence at the selected probability threshold.

\section{Cross-Task Switch Comparison}
\label{sec:results-cross-task-switch-comparison}

The previous sections compare models inside each task. A direct cross-task
comparison requires an additional alignment step, because the tasks do not make
decisions on exactly the same timestamps. The main comparison uses the common
Perfect Switch reference grid and treats missing native decisions as switch
\(0\). This is intentionally conservative: a method is penalized if its task
definition does not produce decisions over part of the reference grid. A second
strict-intersection comparison was also generated, but it is used only as a
diagnostic because it keeps only timestamps where all methods have native
decisions.

\Cref{tab:cross-task-best-results} reports the best method from each task on
the reference-grid comparison. Current-level XGBoost obtains the highest F1 and
perfect precision, but it does so with lower recall than the survival and
long-fade models. XGBoost-AFT is the strongest survival method and is close to
current-level XGBoost in F1 while reaching higher recall. Among autoregressive
models, PatchTST raw gives the best balance. Long-fade detection reaches full
recall, but its precision is lower because it keeps the switch active over a
larger part of grouped fade regions.

\begin{table}[htbp]
  \centering
  \caption{Best switch method from each task on the common reference grid. Missing native decisions are treated as switch \(0\).}
  \label{tab:cross-task-best-results}
  \small
  \setlength{\tabcolsep}{4pt}
  \renewcommand{\arraystretch}{1.15}
  \begin{tabular}{
    @{}
    >{\raggedright\arraybackslash}p{0.25\textwidth}
    >{\raggedright\arraybackslash}p{0.28\textwidth}
    rrrr
    @{}
  }
    \toprule
    Task & Best method
      & \shortstack{Coverage\\\(\uparrow\)}
      & \shortstack{Precision\\\(\uparrow\)}
      & \shortstack{Recall\\\(\uparrow\)}
      & \shortstack{F1\\\(\uparrow\)} \\
    \midrule
    Current-level persistence & XGBoost scalar context & 89.0\% & 1.000 & 0.808 & 0.894 \\
    Survival persistence & XGBoost-AFT & 9.2\% & 0.812 & 0.938 & 0.870 \\
    Autoregressive forecasting & PatchTST raw & 89.2\% & 0.779 & 0.821 & 0.800 \\
    Long-fade detection & TCN classifier & 28.7\% & 0.557 & 1.000 & 0.715 \\
    \bottomrule
  \end{tabular}
\end{table}

\Cref{fig:cross-task-f1-precision,fig:cross-task-recall-balanced-accuracy}
show the same comparison across all evaluated methods: F1 and precision in
the first figure, recall and balanced accuracy in the second. Both use the
same method order, ranked by F1. All values refer to the final switches after
the common island-removal and holding rules; raw model decisions are retained
only as diagnostics.

The F1 and precision panels highlight current-level XGBoost as the most
conservative useful method, with no false-positive switch samples on the
reference grid. PatchTST raw provides the strongest autoregressive balance.
However, precision alone does not show how much reference switch time is
missed, so these results must be read alongside recall.

The recall and balanced-accuracy panels show the other side of this trade-off.
Survival XGBoost-AFT recovers more reference switch time while keeping
precision relatively high. Long-fade detection models reach full or nearly full recall but
introduce extra switch duration, explaining their lower precision in
\cref{fig:cross-task-f1-precision}. The best model therefore depends on the
relative cost of missed coverage and unnecessary activation.

\begin{figure}[H]
  \centering
  \includegraphics[width=\textwidth]{results/cross_task_reference_grid_f1_precision.png}
  \caption{Cross-task F1 (top) and precision (bottom) on the common reference grid. Missing native decisions are mapped to switch \(0\), so the comparison includes both decision quality and task coverage.}
  \label{fig:cross-task-f1-precision}
\end{figure}

\Cref{fig:cross-task-coverage} shows why native decision coverage is a critical
part of the interpretation. Autoregressive and
current-level models operate on most of the reference grid because their
datasets are based on candidate event windows. Long-fade detection is defined
only inside grouped fade events, and survival persistence is defined only while
the sample is inside a detected survival fade event.

\begin{figure}[H]
  \centering
  \includegraphics[width=\textwidth]{results/cross_task_reference_grid_recall_balanced_accuracy.png}
  \caption{Cross-task recall (top) and balanced accuracy (bottom) on the common reference grid, with the same method order and missing-decision policy as \cref{fig:cross-task-f1-precision}.}
  \label{fig:cross-task-recall-balanced-accuracy}
\end{figure}

Consequently, survival
metrics on the reference grid evaluate a smaller native decision domain, even
though the predictions are meaningful inside that domain.
The strict-intersection diagnostic removes this coverage effect by retaining
only the 771 timestamps where all task formulations provide native decisions.
On this reduced grid, current-level XGBoost remains the best conservative
method with precision \(1.000\), recall \(0.819\), and F1 \(0.900\).
XGBoost-AFT obtains precision \(0.814\), recall \(0.938\), and F1 \(0.872\),
and XGBoost raw is the strongest autoregressive method with F1 \(0.866\).

\begin{figure}[H]
  \centering
  \includegraphics[width=1\textwidth]{results/cross_task_reference_grid_native_decision_coverage.png}
  \caption{Native decision coverage for each method on the common Perfect Switch reference grid. Lower coverage means that the corresponding task produces decisions on a smaller subset of reference timestamps.}
  \label{fig:cross-task-coverage}
\end{figure}


The long-fade TCN reaches recall \(1.000\) but only F1 \(0.778\). Therefore,
the main ranking is not only an artifact of missing decisions on the reference
grid: the same qualitative pattern remains when the comparison is restricted
to the shared native decision domain.

\Cref{fig:cross-task-event-f1-heatmap} reports event-level F1 values for the
selected reference events. The heatmap shows that no formulation dominates
uniformly. Some events are handled well by most models, while others expose
clear differences in onset timing, release timing, or excessive switch
activation. This supports the use of both global and event-level metrics:
global F1 summarizes average behavior, whereas event-level views reveal which
fade morphologies remain difficult.

The three all-zero rows corresponding to the current-level shapelet models are
genuine results rather than missing values or a coverage artifact. As discussed
in \Cref{sec:results-current-level-switch}, none of their test predictions
satisfies both the \SI{10.0}{dB} signal gate and the \SI{300}{s} predicted
persistence requirement. Their final switch sequence is therefore zero over
the complete test set. Every selected event containing a positive Perfect
Switch interval consequently has no true-positive shapelet decision and an
event-level F1 of zero. The heatmap makes explicit that this failure is
systematic across events, not confined to a few difficult fade morphologies.

\begin{figure}[H]
  \centering
  \includegraphics[width=\textwidth]{results/cross_task_reference_grid_event_f1_heatmap.png}
  \caption{Event-level F1 heatmap for the cross-task reference-grid comparison. Rows are model-derived switch methods and columns are selected external-holdout events.}
  \label{fig:cross-task-event-f1-heatmap}
\end{figure}

\Cref{fig:cross-task-event-timeline} gives a concrete example on a globally
aligned reference event, \texttt{dataset\_007\_event\_00007}, also used in
\cref{fig:autoregressive-patchtst-switch-example,fig:current-level-xgboost-switch-example}.
This is the figure to use when comparing methods from
different task formulations on exactly the same timestamps. The all-method
timeline makes the precision-recall trade-off visible: conservative models
activate later or for shorter intervals, while high-recall models cover more of
the fade but may extend beyond the Perfect Switch interval.

\begin{figure}[H]
  \centering
  \includegraphics[width=1.0\textwidth]{results/cross_task_event00007_all_methods_switches.png}
  \caption{Cross-task event timeline for the selected multi-interval external-holdout event. The upper panel shows the raw signal and threshold. Each name labels the switch band immediately below it; filled bands show final post-processed decisions and black ticks show raw decisions.}
  \label{fig:cross-task-event-timeline}
\end{figure}

\section{Summary of Main Empirical Findings}
\label{sec:results-main-findings}

The experimental results lead to four main findings.

First, native predictive accuracy and switch quality are related but not
equivalent. In the autoregressive task, the raw GRU models obtain strong
forecast losses, but PatchTST raw gives the best switch F1 because its errors
produce a more useful threshold-crossing pattern. In survival persistence, the
discrete-time TCN has the best integrated Brier score, but XGBoost-AFT gives
the better switch sequence. In long-fade detection, near-perfect AUPRC does not
translate into a precise switch controller because the grouped-event label
overextends the active decision. These cases confirm that final evaluation
must be performed after switch conversion, not only in the native predictive
space.

Second, tabular gradient-boosted models are strong baselines. XGBoost is
competitive in autoregressive forecasting, is clearly the best current-level
persistence model, and XGBoost-AFT is the best survival switch model. This is
important methodologically: more complex neural sequence models are useful, but
they do not automatically dominate well-tuned lag and scalar representations on
this dataset.

Third, the different task formulations express different operating policies.
Current-level XGBoost is conservative and precise; survival XGBoost-AFT is a
strong residual-duration formulation with higher recall; PatchTST raw provides
the best autoregressive switch balance; long-fade detection is a robust
high-recall classifier but tends to overextend switch activity. The
learnable-shapelet current-level models are the clearest negative result: after
the operational conversion they never request a switch. These behaviors are
more informative than a single global ranking.

Fourth, the shared post-processing is essential for a fair operational
comparison. All final switch sequences are evaluated after the same island and
holding rules, and all are compared with the same Perfect Switch construction.
The remaining differences therefore come from the task formulation and model
outputs, not from task-specific switch logic. This also means that the reported
switch metrics should be read as operational metrics of the post-processed
decision sequence, not as metrics of the raw model threshold crossings.

Overall, the most promising formulations are current-level persistence with
XGBoost and survival persistence with XGBoost-AFT, followed by autoregressive
PatchTST raw. Current-level persistence gives the cleanest switch sequence in
the reported reference-grid comparison, but its underlying duration regression
has large tail errors. Survival persistence is conceptually closer to the
operational question because it estimates residual fade duration directly.
PatchTST raw remains valuable because it predicts the future signal trajectory
and can support analyses beyond the binary switch decision. Long-fade detection
is useful as a high-recall event classifier, but less precise as a standalone
switch controller.
